\section{The sieve rows S1--S10 on the actual source}
\label{p2:sec:rows}
This section treats the direct sieve rows on the actual source.
Put $\rho=1+\tau$, $\mu=25/128$. The rows are grouped as
\[
\begin{array}{c|l}
\mathrm{DIRECT}&\mathrm{S1\hbox{-}0\hbox{--}S1\hbox{-}3,S3\hbox{-}A,S4,S5\hbox{-}low,S7,S8\hbox{-}A,S9\hbox{-}A,S10\hbox{-}Ret}\\
\mathrm{S2}&\mathrm{S2\hbox{-}0\hbox{--}S2\hbox{-}3}\quad\text{(whole-$K$ at depth three)}\\
\mathrm{BV}&\mathrm{S5\hbox{-}top,S6}.
\end{array}
\]
Sections~3 and~5 treat S2 and BV.
Conditions C1--C5 are those of Theorem~\ref{thm:distribution-input}.
The internal Lemma~\ref{thm:packing-input} gives the slots.

\subsection{Source replacement and the exact signed coefficient}
This subsection replaces the literal source by the plain one.
\begin{lemma}[Source replacement and coefficients]
\label{lem:source-replacement}
Let $\mathcal P_N=\{p>2:p\nmid N\}$ and
\[
 A_D(v)=\mathbf1_{\mathbb P}(v)
 \mathbf1_{N/\log^{30}N<v<(1-\epsilon_A)N}.
\]
For a squarefree prefix $b$ from $\mathcal P_N$, set
\[
\begin{aligned}
\mathcal A_b&=\{v\in\mathbb P:v<(1-\epsilon_A)N,\ v\nmid N,\ b\mid N-v,
                         \ p^2\nmid N-v\ (p\mid b)\},\\
\mathcal A_b^{\rm pl}&=\{v:A_D(v)=1,\ b\mid N-v\}.
\end{aligned}
\]
For $A\subseteq\mathcal A_b\cup\mathcal A_b^{\rm pl}$, the notation
$S(A,\mathcal P_N,z)$ means
\[
 \#\left\{v\in A:\left(\frac{N-v}{b},P_N(z)\right)=1\right\}.
\]
Thus this sieve tests the cofactor, not the prime $v$ or the integer $N-v$.
If every prefix prime is $\ge z$, the literal sieve is
$S(\mathcal A_b,\mathcal P_N,z)$. Its source comparison is
\[
 S(\mathcal A_b,\mathcal P_N,z)-S(\mathcal A_b^{\rm pl},\mathcal P_N,z)
 =S(\mathcal A_b\setminus\mathcal A_b^{\rm pl},\mathcal P_N,z)
  -S(\mathcal A_b^{\rm pl}\setminus\mathcal A_b,\mathcal P_N,z).
 \tag{2.17a}\label{p2:eq:direct-source-comparison}
\]
The normalized unsifted bounds, in Section~6 units, are
\[
\begin{array}{c|c}
v\le N/\log^{30}N&O_u(L_N^{-28})\\
v\mid N&O_u(L_N^3/N)\\
p^2\mid N-v\text{ for a prefix prime}&O_u(N^{-u/(1+\tau)}L_N^3).
\end{array}
\]
Fix an AA37 sign $\nu$, a row cell and tail labels.
All literal clipped sets and the smooth divisor are kept. Put
\[
 C_c(q)=\sum_{d,(p_j)}\lambda_d\mathbf1_{q=d\prod_jp_j},\qquad |\lambda_d|\le1.
\]
Let $K_T$ be the equal-label tail mask,
and let $K_P$ impose complete equal-label prefix distinctness.
For a complete permutation-invariant block, let $M_c$ be its prefix orbit factorial.
Then
\[
 O_c(q)=\frac{C_c(q)}{M_c}-\Xi_{U,c}(q),\qquad
 \Xi_{U,c}=\frac1{M_c}\sum_{\rm full}\lambda_d(1-K_TK_P)
                  \mathbf1_{q=d\prod_jp_j}.
 \tag{2.17b}\label{p2:eq:direct-c1}
\]
With $\gamma_c=C_c/C_{\rm slot}$, the full remainder is
\[
 R_{{\rm AA},c}=\frac{C_{\rm slot}}{M_c}
       \sum_q\gamma_c(q)\Delta_P(q)
       -\sum_q\Xi_{U,c}(q)\Delta_P(q),
 \qquad \gamma_c=C_c/C_{\rm slot}.
 \tag{2.17c}\label{p2:eq:direct-remainder}
\]
The external multiplier $c\nu$ acts on both terms.
This includes the four S2 children and the whole-$K$ last child.
The same identity holds for nonnegative occurrence weights.
On the common occurrence space, $a_r,b_r$ are the original and plain weights.
The prefix and sifting data are fixed, and the weights are extended by zero. Put
\[
 h_r=\min(a_r,b_r),\qquad
 e_r^+=(a_r-b_r)_+,\qquad e_r^-=(b_r-a_r)_+.
\]
Then $a_r=h_r+e_r^+$, $b_r=h_r+e_r^-$ and $e_r^+e_r^-=0$.
Linearity gives $S_z(a_r)=S_z(b_r)+S_z(e_r^+)-S_z(e_r^-)$,
including the G1 representation multiplicities.
\end{lemma}
\begin{proof}
For $m=(N-v)/b$, reduction modulo a common prime gives
\[
 (m,N)=1\iff v\nmid N,\qquad
 (m,b)=1\iff p^2\nmid N-v\quad(p\mid b).
\]
Here $v\nmid b$, since otherwise $v\mid N$.
These conditions concern the ambient source; cofactor sieving is separate.
Additivity proves \eqref{p2:eq:direct-source-comparison}.
The two differences have the three stated exceptional conditions.
Each point is assigned to the first true condition,
and the resulting unsifted masses are bounded.
These source terms are removed before HB.
A common source point removed by sieving creates no source error.
All later Buchstab and distribution calculations use $A_D$.

For (2.17b) we use $1=K_TK_P+K_T(1-K_P)+(1-K_T)$.
Overlapping repetitions are assigned to the first forbidden-tail condition.
Each modulus of $\Xi_{U,c}$ repeats a prime $\ge L_c^{\sigma^2}$.
Theorem~\ref{thm:repeated-primes} bounds the complete nineteen-row union,
including prefix multiplicities, once.
The smooth divisor lies wholly in slot 1.
Then the finite slot convolution is $C_c/C_{\rm slot}$ on every modulus, repeats included.
Empty rough slots are unit atoms, and each realization and factorial occurs once.

Put $C_A(q)=\sum_{v\equiv N\,(q)}A_D(v)$,
$X_A=\sum_vA_D(v)$, $Y_A(q)=\sum_{(v,q)=1}A_D(v)$ and
\[
 \Delta_0=C_A-X_A/\varphi(q),\qquad \Delta_P=C_A-Y_A/\varphi(q).
\]
Prefixes are $\ge z$ and tails are $<z$.
Totient multiplicativity converts the relative AA37 comparator into $\Delta_0$.
Theorem~\ref{thm:repeated-primes} gives $q<N^{3/5}$ on the full labelled support.
Eventually $v>N/L_N^{30}>N^{3/5}$.
Thus $Y_A=X_A$, so $\Delta_0=\Delta_P$ before HB.
This proves \eqref{p2:eq:direct-remainder}.
\end{proof}

For the actual levels put
\[
\begin{aligned}
v(t)&=\min\{(1+t)/2,(32-34t)/39,57/64-3t/2\},\\
\psi(x,y)&=x\mathbf1_{x\ge y},\\
w(a,b,c)&=\psi\!\left(\min\{(132-68c)/214,(153-96c)/256,1-2b\},(1+a)/2\right),\\
H(a,b,c)&=\max\{v(a),v(b),v(a+b),v(a+b+c),\\
 &w(a,b,c),w(a,c,b),\\
 &\psi(v(a+c),a+2b+c),\\
 &\psi(v(b+c),2a+b+c)\}.
\end{aligned}
 \tag{2.17d}\label{p2:eq:row-level-function}
\]
Let $\widetilde H=H$ when the largest coordinate is $\le\mu$,
and above $\mu$ let $\widetilde H(a,b,c)=v(a)$.
A parent label $t_i$ retains $[N^{t_i},N^{\rho t_i})$, clipped by its row.
At depth $r<3$, complete the ordered triple in $[0,w_z]$, where $z=N^{w_z}$.
Let $\Theta_c$ be the infimum on this closed ordered parent domain.
At depth three the infimum is taken on the three parent intervals.
This infimum is taken before any branch is selected, including crossings at $\mu$.
At depth zero, with $z=N^u$, $\Theta_c=H(u,u,u)=19101/32000$.
Then
\[
 L_c=N^{\ell_c},\qquad \ell_c=\Theta_c-\eta_L-\sum_{i\le r}t_i,
 \qquad D_c=N^{\Theta_c-\eta_L}.
 \tag{2.17e}\label{p2:eq:row-relative-level}
\]
For nonempty DIRECT rows, $t_1<t_*=7109/28000$ and
$509/1000\le\Theta_c\le153/256$. Also
\[
 \Theta_c\ge m(t_1):=\min\{(1+t_1)/2,57/64-3\rho t_1/2\}.
 \tag{2.17f}\label{p2:eq:row-level-floor}
\]
The first branch is active through $t_{\rm dd}=25/[32(1+3\rho)]$.
The required prefix masks are
\[
 t_1+\cdots+t_{j-1}+2t_j<\Theta_c-\eta_L
 \qquad(1\le j\le r).
 \tag{2.17g}\label{p2:eq:direct-prefix-mask}
\]
Their row checks are
\[
\begin{array}{c|l}
\text{rank one}&2t_1<2t_*=7109/14000<509/1000-\eta_L\\
\mathrm{S1\hbox{-}2,S1\hbox{-}3}&t_i<\alpha,\quad4\alpha<509/1000-\eta_L\\
\mathrm{S2\hbox{-}2,S2\hbox{-}3}\ (j=2),\ \mathrm{S8\hbox{-}A}
 &3\beta<509/1000-\eta_L\\
\mathrm{S2\hbox{-}3}\ (j=3)&t_1<\beta<t_{\rm dd}\\
 &t_1+2t_2+(4+2\tau)t_3\le249/250+2\tau\beta<1-2\eta_L\\
\mathrm{S9\hbox{-}A}&t_1<\mu,\ t_2<\beta:\quad t_1+2t_2<m(t_1)-\eta_L\\
\mathrm{S10\hbox{-}Ret}&\Theta_c-u-t_1-t_2\ge(11/5)\alpha\\
 &t_2<\beta<u+(11/5)\alpha-\eta_L.
\end{array}
\]
S2 at rank three uses the first branch of \eqref{p2:eq:row-level-floor}.
S9 substitutes in both branches.
For S10 we use \S\ref{p2:sec:2.10}; its binding prefix-mask gap is
\[
 u+(11/5)\alpha-\eta_L-\beta
 =\frac{12449531}{394490000}>0.03155.
\]
These checks include clipped cells and weak selection endpoints.
AA37 tail labels $D_j=N^{s_j}$ satisfy $\sum_{i<j}s_i+2s_j<\ell_c$.
Together with \eqref{p2:eq:direct-prefix-mask}, prefix/tail separation gives
$\sum_{i<j}\lambda_i+2\lambda_j<\Theta_c-\eta_L$ for merged sorted labels.
Sorted upper endpoints are at most $\rho$ times their labels,
and their mask is strictly below $\rho(\Theta_c-\eta_L)$.
This proves packing geometry, including empty prefixes.
The tested-source and boundary hypotheses are proved below.

For the Section~4.1 representation domains let $R_j$ count ordered representations.
In $R_j^{\rm tag}$ the coordinates are prime to $N$
and $e/(p_1p_2)$ is prime to $p_1p_2$. Put
\[
\begin{aligned}
 A_G&=(2R_{11}+R_{12}+R_{13})\mathbf1_{\epsilon_AN<e<N},\\
 w_{\rm tag}&=(2R_{11}^{\rm tag}+R_{12}^{\rm tag}+R_{13}^{\rm tag})
                    \mathbf1_{\epsilon_AN<e<N},\\
 w_{\rm pl}&=A_G^{\rm cut}=A_G\mathbf1_{e\le N_{\rm cut}},\quad
 N_{\rm cut}=N(1-(\log N)^{1-C_b}),\quad C_b=334,\\
 U_G&=w_{\rm tag}\mathbf1_{e>N_{\rm cut}},\quad
 E_G=(A_G-w_{\rm tag})\mathbf1_{e\le N_{\rm cut}}.
\end{aligned}
\]
Thus $U_G,E_G\ge0$ and $w_{\rm tag}-w_{\rm pl}=U_G-E_G$.
For $\chi_{b,z}(e)=\mathbf1_{b\mid N-e}\mathbf1_{(N-e,P_N(z))=1}$,
\[
 \sum_e w_{\rm tag}\chi_{b,z}=\sum_e w_{\rm pl}\chi_{b,z}+\Delta_G(b,z),
 \qquad |\Delta_G(b,z)|\le\sum_e(U_G+E_G).
\]
Here $U_G$ is the upper cut; $E_G$ counts coordinate divisibility and repetitions.
Section~4 bounds their unsifted masses and local hypotheses.
The source is replaced before Buchstab, and the signed parent correction is retained.
C1--C5 test the untagged discrepancy $\Delta[A_G^{\rm cut}]$ on interior boxes.
Fixing two coordinates at a prime $p$ leaves at most two positive coordinates.
The elementary divisor bound, over all pairs and four source forms, gives
\[
 O\!\left(\sum_{p\ge N^\alpha}\frac N{p^2}(1+\log N)^2\right)
 =O(N^{1-\alpha}(1+\log N)^2).
\]
Division by $\mathcal W\gg N/(\log N)^2$ gives $O(N^{-\alpha}(\log N)^4)$.
This is within the Section~6 source-restoration allowance.

We use the classical BFI hypothesis (A2), \cite{ref2}.
We fix $\mathscr G$ and $B_\gamma>0$.
For every $A>0$ it asks for a uniform $C_A$ such that
\[
 \left|\sum_{\substack{n\equiv\ell\ (k)\\(n,d)=1}}\gamma_n
 -\frac1{\varphi(k)}\sum_{(n,dk)=1}\gamma_n\right|
 \le C_A\|\gamma\|_2T^{1/2}\tau(d)^{B_\gamma}(\log(2T))^{-A}.
\]
Here $\gamma\in\mathscr G$, $\operatorname{supp}\gamma\subseteq[1,T]$, $T\ge2$,
$d,k\ge1$, $\ell\ne0$, $(\ell,k)=1$.
The constant is uniform in the sequence, subinterval and auxiliary coprimality divisor.
For Class II we use the selected elementary factor
and exchange factor names as in the classical theorem.
Theorem~\ref{thm:distribution-input} requires this $\ell^2$ normalization.
The complementary convolution has its separate norm bound.


\subsection{Common direct preparation}
This subsection verifies the local and global hypotheses for the direct families.
\begin{lemma}[Common preparation]
\label{lem:s2-preparation}\label{lem:direct-preparation}
Let $0<\sigma\le10^{-10}$, $\tau=\sigma^9$, $C_b=334$.
Then $C_b\ge\lceil2A+32\rceil$ for $0<A\le(334-32)/2$.
The per-box precision $A'$ is arbitrary; it is distinct from the
prime-source target $A$ in the condition $C_b\ge\lceil2A+32\rceil$.
Each DIRECT and S2 family satisfies H1--H5 and LEN of
\cite[Theorem~3.2]{Z1} on retained Class-II boxes.
For each family $\mathscr F$, separately, H6 and H7 hold with
\[
\begin{gathered}
 E_{\rm log,\mathscr F}\ll_\sigma NL_N^{18-C_b},\quad
 E_{\rm bdry,\mathscr F}\ll_\sigma NL_N^{16-C_b/2},\quad L_N=\log N,\\
 \sum_{j\in\mathcal J_{\mathscr F}}|c_j|\ll_{\sigma,C_b}L_N^{K_*},\qquad
 K_*=6(C_b+1)+2=2012<24C_b+26=8042.
\end{gathered}
\]
Every other nonboundary box is Class I-1 or I-2.
For fixed $A_I>0$, its signed discrepancy before external coefficients is
$O_{\sigma,A_I}(X(\log X)^{-A_I})$.
The entire S2 Class-I error is $o(W_N)$, with
\eqref{p2:eq:s2-normalization} defining $W_N$.
Section~3 treats the strict BV complement separately.
DIRECT and S2 are each counted once.
\end{lemma}
\begin{proof}
Put $L=\log N$, $H=N/L^{30}$, $\Gamma_k=C_k/M_k$,
$P=\sum_{i\le r}t_i$, $\ell=\Theta_c-\eta_L-P\ge\ell_0$.
The source is $A_D(v)=\mathbf1_{\mathbb P}(v)\mathbf1_{H<v<(1-\epsilon_A)N}$.
Source exceptions have been removed.
All prefix, sieve and whole-cell conditions are modulus-side.
The row masks and AA37 bound $\sum s_j<\ell$ give
\[
\begin{gathered}
 q<N^{\rho P+(\rho+\sigma)\ell}
 \le N^{(1+\tau+\sigma)(153/256-\eta_L)}<N^{e_D},\\
 e_D=153/256-\eta_L/2<153/256.
\end{gathered}
\]
For the dyadic $X$-shell covering $v>H$, one has $X>H/2$;
eventually $q\le X^{3/5}$ and $q<H$.
All factors avoid $2N$.
The tested prime exceeds $q$, so its comparators agree before HB.
Let $u_c=L_c^{\sigma^2}\ge N^{\ell_0\sigma^2}$.
Primes below $u_c$ determine the smooth divisor uniquely.
Labelled primes are $\ge u_c$, since $u_c<N^u$.
Thus $u_c^m\le q<N^{3/5}$ gives
\[
 |\Gamma_k(q)|\le m!/M_k\le C_\sigma,\qquad
 m\le\left\lceil\frac3{5\ell_0\sigma^2}\right\rceil.
\]
This includes repeated primes and each numerical realization once.
Three prefixes, $k<\sigma^{-2}$ tails and two signs give $O_\sigma(1)$ blocks.
Whole-cell conditions remove blocks, and smooth divisors remain inside coefficients.

We apply \cite[(B.7)--(B.13)]{Z1} with
$B=30$, $\varepsilon=\epsilon_A$, $a=N\ne v$.
For $j\le3$ we take $s=2j\le6$, $n\in[X,2X)$, $m_i\le2X^{1/3}$.
Factors are $z_0(X)$-rough, including the unit atom.
Put
\[
\begin{gathered}
 \delta_X=(\log X)^{-C_b},\quad J_X=\lceil\log2/\log(1+\delta_X)\rceil,
 \quad r_X=2^{1/J_X},\quad\Delta_X=r_X-1,\\
 \delta_X/3\le\delta_X(1-2\delta_X)\le\Delta_X\le\delta_X.
\end{gathered}
\]
The complete tested-side faces are
\[
 n=H,\quad n=(1-\epsilon_A)N,\quad n=X,2X;\qquad
 m_i=2X^{1/3}\quad(1\le i\le j).
\]
Crossing boxes are removed.
Roughness stays in the coefficients, and whole-$K$ stays modulus-side.
The bound $2X(1+\delta_X)^6<3X$ makes $n\le8X$ redundant.
The cited construction uses these data and the proved coefficient bound,
so it permits prefixed moduli.
Proper powers contribute $O_\sigma(N^{1/2+o(1)})$:
$O(N^{1/2}\log N)$ terms, each $N^{o(1)}$.
The signed HB coefficient is
\[
 c_j(\mathbf v)=(-1)^{j-1}\binom3j
       \prod_{i\le j}\mu_{z_0}(m_i)\log n_1,\qquad |c_j(\mathbf v)|\le3L.
\]

For H6 put $F(n)=d(N-n)+2(1+L)$.
Divisor counting and $\sum_{q\le N^{3/5}}1/\varphi(q)\ll1+L$ give
$|\sum_q\Gamma_k(q)\mathcal S_n(q;N)|\ll_\sigma F(n)$.
Each product fixes its shell by \cite[(B.11)]{Z1}.
We use $\sum_{n\le y}\tau_h(n)\le y(1+\log y)^{h-1}$,
$\tau_6^2\le\tau_{36}$ and $d^4\le\tau_{16}$.
Cauchy--Schwarz gives $(35+15)/2=25$, $(5+25)/2=15$, and
\[
 \sum_{n<N}\tau_6(n)\ll NL^5,\quad
 \sum_{n<N}\tau_6(n)F(n)^2\ll NL^{25},\quad
 \sum_{n<N}\tau_6(n)F(n)\ll NL^{15}.
\]
For $\lambda_B=\sum_i\log V_i$, we have
$\log n,\lambda_B\asymp L$, $0\le\log n-\lambda_B\le6\delta_X$,
$\log X=L+O(\log L)$.
Hence
\[
 |E_{\rm log,\mathscr F}|\ll_\sigma
 L^{-C_b-1}\sum_{n<N}\tau_6(n)F(n)\ll_\sigma NL^{14-C_b}.
\]
This is stronger than stated.
Each tuple has one shell and one box.
For faces we use \cite[Lemma~B.2, (B.14)--(B.16)]{Z1}, with
$2\le s\le6$, $Y_0=N^{1/7}$, $\eta_{\rm dist}\ll L^{-C_b}$.
Source faces are singletons, and the lattices are
$(\log2)\mathbb Z$ and $\log2+(\log2/3)\mathbb Z$.
Active product coordinates exceed $(H/2)^{1/6}$.
At caps, $m_i\asymp X^{1/3}$ and $X>H/2$.
Both exceed $Y_0$ eventually. Thus
\[
\begin{gathered}
 \#\mathcal U\ll NL^{5-C_b}+N^{6/7}L^5\ll NL^{5-C_b},\\
 |E_{\rm bdry,\mathscr F}|\ll_\sigma
 L(\#\mathcal U)^{1/2}(NL^{25})^{1/2}\ll_\sigma NL^{16-C_b/2}.
\end{gathered}
\]
Here $|c_j|/\lambda_B\ll L$.
Finite faces and unique shells introduce no box count.
Three patterns and $O_\sigma(1)$ blocks alter only constants.
This proves H6 separately for both families.

The classification uses \cite[(B.20)--(B.24), (B.27)--(B.28), (B.30)]{Z1}.
We choose $N_0(\sigma)$ so that
\[
 \begin{gathered}
 \omega=\log_N(N/X),\quad\kappa=(\log2+6\log r_X)/L,\\
 \omega,\kappa\le\sigma/202,\quad
 \sigma L\ge2^{130}(C_b+1)(\log L+1).
\end{gathered}
\]
The product and M\"obius caps are exactly (B.21).
The selected variable has the largest eligible lower exponent in $[60\sigma,1/3+14\sigma]$,
with ties broken by index.
The exclusions are $14\sigma>\kappa$, $360\sigma<1-\omega-\kappa$,
$1+42\sigma>1-\omega+\kappa$.
Every other box has one or two large smooth variables.
For the selected Class-II cell $[T,r_XT)$ put
\[
 I_T=(T,r_XT],\quad N_I=\lfloor r_XT\rfloor-\lfloor T\rfloor,\quad
 M_I=M(r_X^r-1),\quad m_0=M(1+r_X^r)/2.
\]
Here $M$ is the complementary lower product, $TM\asymp X$.
Its nonunit-cell count satisfies $1\le r\le5$,
since eventually $n>H>N^{1/3+15\sigma}$.
By \cite[Lemma~B.3, (B.22)--(B.24)]{Z1}, internal endpoints are irrational.
Rough coefficients vanish at powers of two.
Both integer counts differ from $T\Delta_X$ by less than one.
The threshold gives $T\Delta_X>2$ and $T^{2^{-126}}\ge12(\log X)^{C_b}$.
Consequently
\[
\begin{gathered}
 \frac{T}{6(\log X)^{C_b}}\le N_I\le2T(\log X)^{-C_b},\qquad
 M_I\asymp(X/T)(\log X)^{-C_b},\\
 N_I\ge(2T)^{1-2^{-126}},\quad X^{60\sigma}\le T\le X^{1/3+15\sigma},\\
 \log_NM_I\ge2/3-(14+1/202)\sigma-o(1)\ge2/3-14.02\sigma,\\
 M_I\gg X^{2/3-15\sigma}(\log X)^{-C_b}\ge X^{1/2}.
\end{gathered}
\]
The majorant in \cite[Appendix~B, immediately before (B.27)]{Z1}
equals one on $[M,Mr_X^r]$,
with support $[M-M_I/2,Mr_X^r+M_I/2]$.
For a source point $n_0$, enlarged products satisfy
\[
 n_0r_X[1+\tfrac32(r_X^r-1)]\le n_0(1+11\delta_X)<(1-\epsilon_A/2)N.
\]
Thus $N-mn\ge\epsilon_A N/2$, $|mn-N|\le X^3$, $|m_0|\ll X/T$.
Eventually $N\le XL^{31}\le X^{1+\sigma/20}$.
The row masks and Lemma~\ref{thm:packing-input} give exact repeated-prime slots,
including S1-0 and S2-0, and all spectral inequalities.
In particular,
\[
 Q<TX^{-29\sigma},\quad QRS<M_IX^{-89\sigma},\qquad
 \frac{N_I}{QX^{\sigma/3}}\ge
 \frac{X^{(29-1/3)\sigma}}{6(\log X)^{C_b}}\longrightarrow\infty.
\]
For rough M\"obius, unit and normalized-log factors
we use \cite[(3.25), (3.47)--(3.48)]{Z1}.
The complement is bounded by $\tau_5$.
This proves H1--H5 and LEN separately from H6/H7.

For Class I we apply \cite[Lemmas~3.3 and~3.5, \S3.11, (B.30)]{Z1}.
Here $k<\sigma^{-2}$, $|\Gamma_k|\le C_\sigma$, complement $\le\tau_5$.
The nonnegative weights $1,\log v/L$ have bounded variation and supremum.
Predicates are constant and factors are $z_0(X)$-rough.
The discrepancy remains $\mathcal S_{\prod v_i}(q;N)$.
Assign proper powers, LOG, boundary, II, I-1, I-2 in order.
The classification gives
\[
 e_*-e_D-(1-\omega)/10\ge0.9-e_D-0.9\omega-\kappa-300\sigma>0.3023,
\]
so $V\ge qX^{1/10}$.
For I-2 we order $N_1\ge N_2$. Then
\[
 \begin{gathered}
 M\le X^{(240\sigma+\kappa)/(1-\omega)}\le X^{2.401\cdot10^{-8}},\\
 N_2\ge X^{1/3},\quad N_1N_2M\le2^{30}X,\quad q<X^{0.5977}.
\end{gathered}
\]
Larger $N_0$ absorbs fixed dyadic factors into $\kappa$.
Thus \cite[(3.100)--(3.101)]{Z1} hold, and (3.105) and (3.111) apply.
For H7, at most six coordinates and three patterns have
$O(\log\log N)\le O(L)$ shells.
Each coordinate has $O(L)$ dyadic intervals and $J_X=O(L^{C_b})$ cells.
Labels, signs, binomial factors and slot factorials cost $O_\sigma(1)$.
The sole logarithmic scalar is $L/\lambda_B=O(1)$.
Hence
\[
 \sum_{j\in\mathcal J_{\mathscr F}}|c_j|
 \ll_{\sigma,C_b}L^{6(C_b+1)+1}\le C_{\sigma,C_b}L^{K_*}.
\]
The extra unit covers the I-2 small-product split.
Realizations and smooth divisors remain inside coefficients.
PF and the selected boundary region form one family.
With $A_I=K_*+6$, the Class-I total is $O_\sigma(NL^{-6})$.
Since $W_N\gg NL^{-2}$, its normalized bound is $O_\sigma(L^{-4})$.
Normalized LOG and boundary errors are
$O_\sigma(L^{20-C_b}+L^{18-C_b/2})=o(1)$.
\end{proof}


\subsection{Four-term Buchstab and the S1 comparison}
\label{p2:sec:2.1}\label{p2:sec:2.2}
This subsection proves the Buchstab chains and bounds the cutoff comparison.
\begin{lemma}[Four-term chains]
Set $S_{\rm pl}(b,z)=S(\mathcal A_b^{\rm pl},\mathcal P_N,z)$.
The plain-source identity is
\[
\begin{aligned}
3S_{\rm pl}(1,N^\alpha)
={}&3S_{\rm pl}(1,N^u)
 -3\sum_{u\le r_1<\alpha}S_{\rm pl}(p_1,N^u)\\
&+3\sum_{u\le r_2<r_1<\alpha}S_{\rm pl}(p_1p_2,N^u)
 -3\sum_{u\le r_3<r_2<r_1<\alpha}S_{\rm pl}(p_1p_2p_3,p_3).
\end{aligned}\tag{2.19}\label{p2:eq:2.19}
\]
Here $r_i=\log p_i/\log N$ and all $p_i\in\mathcal P_N$.
The S2 identity (3.2) replaces $\alpha$ by $\beta$ and divides by $3$.
Both use $L_c=N^{\Theta_c-\eta_L-P}$, $P=\sum_{i\le r}t_i$, $A_R=c/2$.
Their parameters are
\begin{center}
\begin{tabular}{@{}cllll@{}}\toprule
Row & $c$; direction & Cutoff & Floor for $\ell_c$ & $M_c$\\\midrule
S1-0 & $3$; lower $f$ & $N^u$ & $19101/32000-\eta_L$ & $1$\\
S1-1 & $-3$; upper $F$ & $N^u$ & $5243851/11490000$ & $1$\\
S1-2 & $3$; lower $f$ & $N^u$ & $1414617/3830000$ & $1,2$\\
S1-3 & $-3$; upper $F$ & $E_3=N^{t_3}$ & $3743851/11490000$ & $1,2,6$\\
S2-0 & $1$; lower $f$ & $N^u$ & $95489/160000$ & $1$\\
S2-1 & $-1$; upper $F$ & $N^u$ & $1294691/3090000$ & $1$\\
S2-2 & $1$; lower $f$ & $N^u$ & $264897/1030000$ & $1,2$\\
S2-3 & $-1$; upper $F$ & $E_3=N^{t_3}$ & whole-$K$: \S\ref{sec:s2-3} & $1,2,6$\\\bottomrule
\end{tabular}
\end{center}
The S1 functional has signs $+,-,+,-$ and densities
$f/u,-F/u,+f/u,-\Phi$. Appendix~A.2 gives
\[
 12.9020516<\mathcal S_1<12.9023390,\qquad\tfrac34\mathcal S_1>9.6765387.
\]
Its cutoff allowance is $62500\tau D$.
The directed prefix comparison costs zero and has multiplicity one.
The true one-prefix conversion error is zero.
The alternative $-f$ repair costs less than $10^{-390}D$.
\end{lemma}
\begin{proof}
For S1, \eqref{p2:eq:direct-source-comparison} gives
$3S_1=3S(\mathcal A_1^{\rm pl},\mathcal P_N,N^\alpha)+3\Delta_1^{\rm src}$.
For S2 we use Section~\ref{sec:s2-1}:
$S_{2,\rm raw}=S(\mathcal A_{2,1},\mathcal P_N,N^\beta)$.
Both parents have prefix $1$, $v<(1-\epsilon_A)N$, $v\nmid N$.
The corrections are low-prime mass minus exceptional $v\mid N$ mass.
Prefix-square restoration is empty, and its nonnegative allowance $O_u(N^{-u/(1+\tau)}L_N^3)$ is retained.
These unsifted bounds are included once.
Descendants use only $A_D(v)$, $b\mid N-v$ and cofactor sieving.
Buchstab's identity holds for every sequence, without squarefree or coprime
restrictions \cite{Buchstab1937}.
Three iterations give (2.19).
The coefficients are $(3,-3,3,-3)$ and $(1,-1,1,-1)$.
In subsequent sieve calculations we write $S=S_{\rm pl}$.
We apply \cite[Theorem~4]{ref5} in the order lower, upper, lower, upper.
For S1-0, $\ell-2u=94849/160000>0$.
For S1-1 and S1-2, respectively, $s>228.19$ and $s>184.67$.
Equal-label pairs retain $[N^u,N^\alpha)$ or $[N^u,N^\beta)$.
The full summand, including $L_c$, depends on unordered products and labels.
Thus \eqref{p2:eq:direct-c1} gives $M_c=2$ there, and $1$ otherwise.
Its diagonal is assigned to prefix squares.

We fix $E_3$ before the structured sieve.
Since $p_3\ge E_3$,
\[
 -3S(p_1p_2p_3,p_3)\ge-3S(p_1p_2p_3,E_3).
 \tag{2.20}\label{p2:eq:2.20}
\]
For S1-3, $\ell-t_3\ge2743851/11490000>0$; the upper sieve is legal.
Fixing $E_3$ makes the complete tail domain permutation-invariant.
The factorials are $1,2,6$.
Formula \eqref{p2:eq:direct-remainder} has scalar $-3\nu C_{\rm slot}/M_c$.
Using actual $p_3$ before this comparison would lose invariance.
All four terms have fixed modulus-side labels.
Lemma~\ref{lem:s2-preparation} supplies their boxes.
The equal-prefix terms have disjoint prefix diagonals.
Theorem~\ref{thm:repeated-primes} bounds their union;
Theorem~\ref{thm:sieve-mass} bounds the AA37 main error.

For the S1-3 comparison put
\[
\begin{aligned}
\Theta_3^{\rm saved}(a,b,c)=\max\{&v(a)-u,v(b)-u,v(a+b)-u,v(a+b+c)-u,\\
 &\psi(v(a+c)-u,a+2b+c),\psi(v(b+c)-u,2a+b+c)\}.
\end{aligned}
\]
\hypertarget{p2:def:six-branch-saved}{}
Here both $w$ branches are deleted.
Every retained $v$ is reduced by $u$, including the guarded first inputs.
Let $\Theta_3^{\rm saved}[\tau\text{-bin}(r)]$ be its infimum
on the full ordered comparison box $[r,\rho r]$.
This box differs from the inherited source cell. Put
\[
\begin{gathered}
 G_\ell(w)=F(\ell/w)/w,\quad
 \ell_p(r)=\Theta_3^{\rm saved}(r)-r_1-r_2-r_3,\\
 \Phi(r)=F\!\left(\frac{\Theta_3^{\rm saved}[\tau\text{-bin}(r)]-\sum r_i}
 {\rho r_3}\right)/r_3.
\end{gathered}
\]
On the full parent, $u\le a,b,c\le\rho\alpha$ and $2\rho\alpha<25/128$.
Thus $v(a+c),v(b+c)\ge1/2$.
Both thresholds are $\le4\rho\alpha<1/2-u$.
The original and saved guards are active throughout.
Adding $u$ restores six corresponding $H$ branches,
each with slopes of absolute value $\le3/2$.
Their maximum is $3/2$-Lipschitz in coordinate $\ell^1$ distance.
For $x\in[r,\rho r]$, $\sum_i|x_i-r_i|\le\tau\sum_i r_i\le3\alpha\tau$.
Therefore
\[
 \Theta_3^{\rm saved}+u\le H,\qquad
 \Theta_c\ge\Theta_3^{\rm saved}(r)+u-(9/2)\alpha\tau
              >\Theta_3^{\rm saved}(r)+\eta_L.
\]
Branchwise, $(v-u)(\sum_{i\in A}r_i)-\sum_i r_i$ is coordinatewise nonincreasing.
The inherited lower labels give the same residual comparison $\ell_p(r)\le\ell_c$.
The comparison-bin infimum is at most its point value.
Decreasing $F$ gives $F(\ell_p(r)/(\rho r_3))/r_3\le\Phi(r)$.
For $H_F(s)=sF(s)$, the delay equations give
$H_F'=0$ on $[1,3]$ and $H_F'=f(s-1)\in[0,1]$ above $3$.
Hence
\[
 G_{\ell_c}(t_3)\le G_{\ell_c}(r_3)+\tau/r_3
 \le(1+\tau)G_{\ell_c}(r_3)\le(1+\tau)\Phi(r).
 \tag{2.21}\label{p2:eq:2.21}
\]
For S1-3, $t_i<\alpha$.
On the full parent, $a+c\le2\rho\alpha<t_{\rm dd}<25/128$,
so $v(a+c)=(1+a+c)/2$.
Also $a+2b+c\le4\rho\alpha<1/2$, so its $\psi$ branch is active.
Since $a\ge t_1$ and $c\ge t_3$,
\[
 \Theta_c\ge\frac12+\frac{t_1+t_3}{2},\qquad
 \ell_c\ge\frac12-t_2-\frac{t_1+t_3}{2}-\eta_L
             \ge\frac12-2\alpha-\eta_L.
\]
Since $\eta_L<u$ and $r_3\le\rho\alpha$,
differentiation is legal:
$\ell_c/r_3>(1/2-u-2\alpha)/(\rho\alpha)=186101/(50000\rho)>1$.
The comparison-box guards also remain active.
Its numerator loses at most $(9/2)\alpha\tau$, and the argument stays $>3.7$.
Thus both monotonicities apply, using the actual bin infimum.
Integration over the same literal clipped partition gives
\[
 E_{\rm cut}\le\tfrac34(2/e^\gamma)\tau I^*<1821.2\tau\,D
 \le62500\tau\,D.
 \tag{2.22}\label{p2:eq:2.22}
\]
Here Appendix~A.2 proves $I^*\le I_0+B_J<2161.074150710003$,
uniformly for $0\le\tau\le10^{-90}$.
Reproduction requires the interval-arithmetic source, dependencies and directed output;
see Appendix~A's data-availability statement.
The allowance $62500\tau D$ includes (2.22),
and Section~6.3.5(c) includes the prefix comparison once.

The one-prefix profile is $-F$, already used by the comparison functional.
For the alternative $-f$ we write $-F=-f-(F-f)$.
The delay equations give $s(F-f)=2e^\gamma\varrho(s-1)$,
where $\varrho$ is Dickman's function.
Here $s\ge s_*=64441/309$, $v_*=s_*-1$.
We iterate $\varrho(v)\le\varrho(v-1)/v$ exactly $207$ times;
the factor is $4/u=2000$.
Since $\log(\alpha/u)\le4$, $\log(\beta/u)\le5$,
\[
 E_{f\to F}\le
 \frac{2000(\tfrac34\cdot4+\tfrac14\cdot5)}
 {s_*\prod_{j=0}^{206}(v_*-j)}<10^{-390.77}<10^{-390}.
 \tag{2.23}\label{p2:eq:2.23}
\]
The S2 part belongs to its atomic endpoint.
Section~6 subtracts the residual $10^{-390}D$ once.
These are the true functions $f,F$.
Equations (2.19)--(2.23) prove the source/profile comparison.
Section~6.3.5(b) uses the single joint S1 bound.

The S2 last child retains the invariant condition (3.3), $h\le K$.
This includes PF and all remaining selected blocks.
Here $E_3$ is fixed before C1, and the scalar is $-\nu C_{\rm slot}/M_c$.
The actual-$p_3$ comparison is kept separate.
The strict complement uses actual $p_3$ and BV.
A PF-only coefficient would omit selected source blocks.
Since $1-2\eta_L-249/250=19/5000$, $K<1-2\eta_L$.
Lemma~\ref{thm:packing-input} uses the literal whole-$K$ infimum.
Section~\ref{sec:s2} checks the local source hypotheses and Class I.
The $X^{1-c_\sigma}$ term of Theorem~\ref{thm:distribution-input} is kept.
The strict BV complement has no prepared Class-II family.
Sections~\ref{sec:s2-3}--\ref{sec:s2-7} prove $d_{T8}$.
They retain whole-$K$ norms, spectral parameters, BV coefficients and the eight-branch functional.
Retain $275009680\tau D$, one structured-sieve error share and tail allowance $D/20000$.
Its true one-prefix profile is $-F$.
Section~6.3.5(b) subtracts these errors once.
\end{proof}


\subsection{Signed fixed-cutoff profiles}
\label{p2:sec:2.3}\label{p2:sec:2.4}\label{p2:sec:2.5}\label{p2:sec:2.6}\label{p2:sec:2.7}\label{p2:sec:2.8}\label{p2:sec:2.9}
This subsection treats the fixed-cutoff rows.
\begin{lemma}[Fixed-cutoff rows]
The six fixed-cutoff rows and S10-Ret have the following data.
In all tables $P=\sum_{i\le r}t_i$, $L_c=N^{\Theta_c-\eta_L-P}$.
The source is $A_D$, with shift $N$.
Write $\mathcal E$ for the common errors listed below.
Let $U_j,L_j$ be the selected raw-profile upper and lower endpoints
of Appendix~A.3, respectively.
Their directed $D$ contributions appear in the last table.
The relative AA37 parameter is $A_R=c/2$.
Each full remainder is \eqref{p2:eq:direct-remainder}, with scalar $c\nu C_{\rm slot}/M_c$.
\begin{center}
\small
\begin{tabularx}{\textwidth}{@{}lXX@{}}\toprule
Row & Exact exponent domain & Parent/anatomy/selection; depth\\\midrule
S3-A & $\alpha\le r_1<\mu$ & Inherited A/B at $t_{\rm dd}$; $1$\\
S4 & $\mu\le r_1<1/4$ & Inherited A/B at $t_{\rm dd}$; $1$\\
S5-low & $1/4\le r_1<t_*$ & B, clipped at $N^{t_*}$; $1$\\
S7 & $\mu\le r_1<\lambda$ & Inherited A/B at $t_{\rm dd}$; $1$\\
S8-A & $\alpha\le r_2<r_1<\beta$ & Small-coordinate anatomy; $2$\\
S9-A & $\beta\le r_1<\mu$, $\alpha\le r_2<\beta$ & Inherited A/B; $2$\\
S10-Ret & $\mu\le r_1<\lambda$,\newline
$\alpha\le r_2<\beta$,\newline $r_2<B_0(r_1)$
 & Whole parent cells, $t_2\le\operatorname{sel}(t_1)$; $2$\\\bottomrule
\end{tabularx}
\end{center}
\begin{center}
\small
\begin{tabular}{@{}lccll@{}}\toprule
Row & $c$; $D$ multiplier & Direction & Cutoff & Relative level\\\midrule
S3-A & $-2$; $-1/2$ & upper $F$ & $N^\alpha$ & $L_c$\\
S4 & $-1$; $-1/4$ & upper $F$ & $N^\alpha$ & $L_c$\\
S5-low & $-1$; $-1/4$ & upper $F$ & $N^\alpha$ & $L_c$\\
S7 & $-1$; $-1/4$ & upper $F$ & $N^\alpha$ & $L_c$\\
S8-A & $1$; $1/4$ & lower $f$ & $N^\alpha$ & $L_c$\\
S9-A & $1$; $1/4$ & lower $f$ & $N^\alpha$ & $L_c$\\
S10-Ret & $1$; $1/4$ & lower $f$ & $N^\alpha$ & $L_c$\\\bottomrule
\end{tabular}
\end{center}
\begin{table}[htbp]
\caption{Direct floors, cutoff gaps and prefix orbits.}\label{tab:direct-parameters}
\small
\begin{tabularx}{\textwidth}{@{}lXXc@{}}\toprule
Row & Lower floor for $\ell_c$ & Cutoff gap; ratio $s$ & $M_c$\\\midrule
S3-A & $64359/160000$ (A), $>0.402243$ (B)
 & $\ell-\alpha>0.315211$; $s>4.6217$ & $1$\\
S4 & $64359/160000$ (A), $\ge0.265524$ (B)
 & $\ell-\alpha>0.178492$; $s>3.0508$ & $1$\\
S5-low & $>0.255792$
 & $\ell-\alpha>0.168760$; $s>2.9390$ & $1$\\
S7 & $64359/160000$ (A), $>0.309274$ (B)
 & $\ell-\alpha>0.222242$; $s>3.5535$ & $1$\\
S8-A & $264897/1030000$
 & $\ell-2\alpha\ge98366653/1183470000$; $s>2.9550$ & $1,2$\\
S9-A & $11886931/49440000$ (A), $>0.240430$ (B)
 & $\ell-2\alpha>0.066366$; $s>2.7625$ & $1$\\
S10-Ret & $11886931/49440000$ (A), $2221831/11490000$ (B)
 & $\ell-2\alpha\ge221831/11490000$; $s>2.2218$ & $1$\\\bottomrule
\end{tabularx}
\end{table}
\begin{center}
\begin{tabularx}{\textwidth}{@{}lXX@{}}\toprule
Row & Saved profile; directed comparison cost & Selected output; estimates outside that endpoint\\\midrule
S3-A & $-\mathcal I_{3,A}/2$; zero & $-U_{3,A}/2$; $\mathcal E$\\
S4 & $-\mathcal I_4/4$; zero & $-U_4/4$; $\mathcal E$\\
S5-low & $-\mathcal I_{5,\rm low}/4$; zero & $-U_{5,\rm low}/4$; $\mathcal E$\\
S7 & $-\mathcal I_7/4$; zero & $-U_7/4$; $\mathcal E$\\
S8-A & $\mathcal I_{8,A}/4$; zero & $L_{8,A}/4$; $\mathcal E$, prefix and tail squares\\
S9-A & $\mathcal I_{9,A}/4$; zero & $L_{9,A}/4$; $\mathcal E$, tail squares only\\
S10-Ret & $\mathcal I_{10}/4$; $7.160\tau$
 & $L_{10}/4$; $\mathcal E$, tail squares only.
 Selected raw interval:
 $\left[\begin{gathered}0.858847455777940,\\0.867486588163666\end{gathered}\right]$\\\bottomrule
\end{tabularx}
\end{center}
S8-A and S9-A have zero cell-selection cost.
S10 has zero profile and S2-tail error.
\end{lemma}
\begin{proof}
The proof applies \cite[Theorem~4]{ref5} and \cite[Theorem~2.6, (2.30)--(2.33b)]{Z1}.
Table~\ref{tab:direct-parameters} gives $\ell_c\ge\ell_0$,
$z\le L_c$ for negative rows and $z^2\le L_c$ for positive rows.
Prefixes are $\ge z$ and tails are $<z$.
Every predicate is modulus-label dependent.
Lemma~\ref{thm:packing-input} and Lemma~\ref{lem:s2-preparation}
give the slots and local/complementary bounds.
Section~6 includes each of the following errors once:
\[
\begin{array}{c}
\text{source restoration, low cut, repeated primes, logarithm freezing, boundary,}\\
\text{Class I, Class II, roughness, structured sieve, prime mass}.
\end{array}
\]
These form $\mathcal E$; its squareful terms have the stated row restrictions.

For one uniform comparison put
\[
\begin{gathered}
 g_3(r)=v(r)-u,\quad
 g_8(r_1,r_2)=\max\{v(r_1),v(r_2),v(r_1+r_2)\},\quad g_9(r)=v(r_1),\\
 g(r)=57/64-u-(5/2)r,\qquad u=1/500.
\end{gathered}
\]
We call the absolute S3-A/S8-A/S9-A profiles Option-A.
Every missing completion contains the corresponding unthresholded branches.
Lipschitz control gives
\[
 \ell_c\ge g_3(r)-r+19/10000-3\tau,\qquad
 \ell_c-[g_8(r)-u-r_1-r_2]\ge19/10000-3\tau>0.
\]
The S4/S7 A/B levels give the same buffer.
For S5-low we use $\Theta_c\ge57/64-3\rho t_1/2$ in B.
On S9-A, $4\beta+3\alpha>25/32$ gives $g_8=g_9$.
Decreasing $F$ and increasing $f$ give all six signed lower comparisons.
Inherited labels are kept through clipping and before branch selection.
For S3, S4, S7 we split at $t_{\rm dd}$.
For S9 we split the inherited anatomy.
Retain both literal two-prefix infima.
S4's clipping at $\mu$ retains both A/B bounds.
Section~6.3 sums its two remainder children once.
S4 and S7 overlap in the initial signed inequality;
their HB partitions remain separate and disjoint for each term.

Put $\mathcal W=4(1-\varepsilon_A)U_C$.
In the displays below, $D_j$ is the signed source contribution divided by
$\mathcal W$, and $\mathcal E_j$ is that row's allocated share of the
common error $\mathcal E$. Section~6 charges these shares once.
The signed sources and their normalized bounds are
\begin{gather}
 \mathcal W D_{3,A}:=-2\sum_{\alpha\le r_1<\mu}S(p_1,N^\alpha).
 \tag{2.24}\label{p2:eq:2.24}\\
 D_{3,A}\ge-\frac12\mathcal I_{3,A}-\mathcal E_{3,A}-o_\sigma(1),\qquad
 \mathcal I_{3,A}=\frac{2}{e^\gamma\alpha}
 \int_\alpha^\mu F\!\left(\frac{v(r)-u-r}{\alpha}\right)\frac{dr}{r}.
 \tag{2.25}\label{p2:eq:2.25}
\end{gather}
\begin{gather}
 \mathcal W D_4:=-\sum_{\mu\le r_1<1/4}S(p_1,N^\alpha).
 \tag{2.26}\label{p2:eq:2.26}\\
 D_4\ge-\tfrac14\mathcal I_4-\mathcal E_4-o_\sigma(1),\qquad
 \mathcal I_4=\frac{2}{e^\gamma\alpha}
 \int_\mu^{1/4}F\!\left(\frac{57/64-u-(5/2)r}{\alpha}\right)\frac{dr}{r}.
 \tag{2.27}\label{p2:eq:2.27}
\end{gather}
For S5 set $t_*=7109/28000$, $\tau_S=57/224$.
Before the structured sieve or HB we split
\[
 -\sum_{1/4\le r_1<\tau_S}S(p_1,N^\alpha)
 =-\sum_{1/4\le r_1<t_*}S(p_1,N^\alpha)
  -\sum_{t_*\le r_1<\tau_S}S(p_1,N^\alpha).
 \tag{2.28}\label{p2:eq:2.28}
\]
The first term is S5-low and the second is S5-top.
The endpoint $t_*$ belongs to BV.
Clipping at $N^{t_*}$ retains parent labels.
S5-low uses B and
\begin{gather}
 \mathcal W D_{5,\rm low}:=-\sum_{1/4\le r_1<t_*}S(p_1,N^\alpha),\notag\\
 D_{5,\rm low}\ge-\tfrac14\mathcal I_{5,\rm low}
                 -\mathcal E_{5,\rm low}-o_\sigma(1),\notag\\
 \mathcal I_{5,\rm low}=\frac{2}{e^\gamma\alpha}
 \int_{1/4}^{t_*}F\!\left(\frac{57/64-u-(5/2)r}{\alpha}\right)\frac{dr}{r}.
 \tag{2.29}\label{p2:eq:2.29}
\end{gather}
For S7,
\begin{gather}
 \mathcal W D_7:=-\sum_{\mu\le r_1<\lambda}S(p_1,N^\alpha).
 \tag{2.30}\label{p2:eq:2.30}\\
 D_7\ge-\tfrac14\mathcal I_7-\mathcal E_7-o_\sigma(1),\qquad
 \mathcal I_7=\frac{2}{e^\gamma\alpha}\int_\mu^\lambda
 F\!\left(\frac{57/64-u-(5/2)r}{\alpha}\right)\frac{dr}{r}.
 \tag{2.31}\label{p2:eq:2.31}
\end{gather}
Its profile loss is zero.
Appendix~A.3 proves its outward upper endpoint,
which Section~6.4 uses at the fixed positive parameter.
Equation (6.5b) uses the S3 comparison.

For the two-prefix rows,
\begin{gather}
 \mathcal W D_{8,A}:=+\sum_{\alpha\le r_2<r_1<\beta}S(p_1p_2,N^\alpha).
 \tag{2.32}\label{p2:eq:2.32}\\
 D_{8,A}\ge\tfrac14\mathcal I_{8,A}-\mathcal E_{8,A}-o_\sigma(1),\notag\\
 \begin{aligned}
 \mathcal I_{8,A}
   &=\frac{2}{e^\gamma\alpha}
     \iint_{\alpha\le r_2<r_1<\beta}
     f\!\left(\frac{g_8(r)-u-r_1-r_2}{\alpha}\right)\\
   &\qquad\qquad\times\frac{dr_1dr_2}{r_1r_2}.
 \end{aligned}
 \tag{2.33}\label{p2:eq:2.33}
\end{gather}
\begin{gather}
 \mathcal W D_9:=+\sum_{\beta\le r_1<\mu\atop\alpha\le r_2<\beta}
 S(p_1p_2,N^\alpha).
 \tag{2.34}\label{p2:eq:2.34}\\
 D_9\ge\tfrac14\mathcal I_{9,A}-\mathcal E_9-o_\sigma(1),\qquad
 \mathcal I_{9,A}=\frac{2}{e^\gamma\alpha}
 \int_\beta^\mu\int_\alpha^\beta
 f\!\left(\frac{v(r_1)-u-r_1-r_2}{\alpha}\right)
 \frac{dr_2dr_1}{r_1r_2}.
 \tag{2.35}\label{p2:eq:2.35}
\end{gather}
Only S8-A has equal-set prefix orbits.
Its equal-label sets are identically clipped to $[N^\alpha,N^\beta)$.
The full summand and level are symmetric before slotting.
Thus $M_c=2$ there and $1$ elsewhere.
Formula \eqref{p2:eq:direct-c1} removes strict triangular ordering.
Its prefix-square part is disjoint from the forbidden-tail union.
The range lies in the small-coordinate anatomy, and both squareful errors are included.
Its outward lower endpoint enters (6.4) once, with coefficient $+1/4$.
S9's literal prefix sets are disjoint, including $N^\beta$.
Hence $M_c=1$, the prefix diagonal is empty,
and both normalization differences vanish by \eqref{p2:eq:direct-remainder}.
Only its tail-squareful error occurs.
Its lower endpoint has coefficient $+1/4$ and multiplicity one;
see Sections~6.3.3 and~6.4.

S5-top and S6 use BV separately.
For S5-top the total level is $N^{1/2}L_N^{-B}$,
with relative level obtained by dividing by the actual prefix.
Here $\kappa=\kappa_{\rm BV}=0$, cutoff $N^\alpha$, profile $-F$.
Section~5 proves its full coefficient and error estimate.
Its prepared Class-II, packing and Class-I contributions are zero,
and its main term enters separately.
Section~6.3 includes S5-low errors once.

\end{proof}
\subsection*{The separate BV row S6}
This subsection bounds the separate BV row.
Its exact source is
\[
 \mathcal W D_6:=-\sum_{\tau_S\le r_1<1/3}S(p_1,N^\alpha),\qquad
 \tau_S=57/224,\quad c=-1,\quad A_R=-1/2.
 \tag{2.36}\label{p2:eq:2.36}
\]
Section~5.3 takes $\kappa_{\rm BV}=10^{-6}<61/766=1/6-\alpha$.
The margin $(1/6-10^{-6}-\alpha)/\alpha$ absorbs the logarithmic loss uniformly.
Put $Q_{\rm BV}=N^{1/2-\kappa_{\rm BV}}L_N^{-B}$,
$D_p=Q_{\rm BV}/p$, with cutoff $N^\alpha$.
Then
\[
 \frac{\log D_p}{\alpha\log N}
 =\frac{1/2-\kappa_{\rm BV}-r_1}{\alpha}
  -\frac{B\log\log N}{\alpha\log N}>1
\]
eventually, uniformly for $r_1<1/3$.
Thus the upper sieve $F$ is legal, and
\[
 \mathcal I_6(\kappa_{\rm BV})=
 \frac{2}{e^\gamma\alpha}\int_{\tau_S}^{1/3}
 F\!\left(\frac{1/2-\kappa_{\rm BV}-r}{\alpha}\right)\frac{dr}{r},
 \qquad D_6\ge-\mathcal I_6(\kappa_{\rm BV})/4-o(1).
 \tag{2.37}\label{p2:eq:2.37}
\]
The BV error includes comparison with the $1/2$-level profile.
Section~5 gives the full weighted BV remainder, prefix multiplicities and endpoint direction.
This row uses the classical proof, without structured slot normalization or C1--C5.
Its BV error is subtracted once, together with its main term.


\subsection{The retained curved S10 row}
\label{p2:sec:2.10}
This subsection bounds the retained curved row.
\begin{lemma}[Retained curved profile]
Put
\[
\begin{aligned}
 C_0&=57/64-u-(11/5)\alpha,\quad
 B_0(r)=\min\{\beta,\max\{\alpha,C_0-5r/2\}\},\\
 B_\tau(t)&=\min\{\beta,\max\{\alpha,C_0-(5/2+3\tau/2)t\}\},\\
 R(t)&=\min\{\rho t,\lambda\},\qquad
 \operatorname{sel}(t)=\min\{B_\tau(t),B_0(R(t))/\rho\}.
\end{aligned}
\]
The actual source is
\[
 \mathcal W D_{10}:=+\sum_{\mu\le r_1<\lambda\atop\alpha\le r_2<\beta,\ r_2<B_0(r_1)}
 S(p_1p_2,N^\alpha).
 \tag{2.38}\label{p2:eq:2.38}
\]
Ret consists of whole parent-label cells with $t_2\le\operatorname{sel}(t_1)$.
For the curved source-trimmed target
\[
 \mathcal I_{10}=\frac{2}{e^\gamma\alpha}
 \int_\mu^\lambda\int_\alpha^{B_0(r_1)}
 f\!\left(\frac{57/64-u-(5/2)r_1-r_2}{\alpha}\right)
 \frac{dr_2dr_1}{r_1r_2}.
 \tag{2.39}\label{p2:eq:2.39}
\]
we have
\[
 D_{10}\ge\mathcal I_{10}/4-7.160\tau-\mathcal E_{10}-o_\sigma(1).
 \tag{2.40}\label{p2:eq:2.40}
\]
The profile and S2-tail errors are zero.
\end{lemma}
\begin{proof}
On clipped support,
$r_2<\rho t_2\le B_0(R(t_1))\le B_0(r_1)$.
Thus $\mathrm{Src}=\mathrm{Ret}\sqcup\mathrm{Drop}$ at finite $N$.
The coefficient $+1$ permits discarding Drop before estimating Ret.
The literal prefixes $[N^\mu,N^\lambda)$, $[N^\alpha,N^\beta)$ are disjoint.
Hence $M_c=1$ and the prefix diagonal is empty.
Selection is constant on whole labels.
For anatomy A, $t_2\le\beta$ gives residual $13475869/473388000>0$.
Table~\ref{tab:direct-parameters} verifies the lower cutoff, including the crossing band.
The tested source stays $A_D$, with shift $N$.
On Ret,
\[
 \alpha(s_{\rm live}-s_{\rm saved})\ge
 (u-\eta_L)-\tau((5/2)\lambda+\beta)>0.
\]
Increasing $f$ gives the favorable direction.
For every discarded actual point,
\[
 t_2>\operatorname{sel}(t_1),\quad r_2\ge t_2,\quad
 r_2<B_0(r_1)\le B_0(t_1).
\]
The larger of the two selection deficits bounds the strip.
The projection $P$ onto $[\alpha,\beta]$ is $1$-Lipschitz, so $B_0$ is $5/2$-Lipschitz.
The two deficits satisfy
\[
\begin{aligned}
 B_0(t)-B_\tau(t)&\le(3/2)t\tau,\\
 B_0(t)-B_0(R(t))/\rho
 &\le(5/2)(R(t)-t)+\beta(1-1/\rho)\\
 &\le(5\lambda/2+\beta)\tau.
\end{aligned}
\]
The projection estimate covers the $\beta$ cap, affine middle and $\alpha$ floor.
The two displayed inequalities cover both branches of the minimum.
Also $R(t)-t\le\tau\lambda$, including the clipped case $R(t)=\lambda$,
and $1-1/(1+\tau)\le\tau$.
Therefore
\[
 \operatorname{width}\le(5\lambda/2+\beta)\tau=\frac{36737}{49440}\tau.
\]
This includes endpoint cells: $t_1\le r_1<\lambda$, $t_2\le r_2$.
The inherited labels at $\mu,\lambda$ are retained.
On the whole canonical source, $11/5\le s<11519839/3200000<4$.
Here $f(s)=2e^\gamma\log(s-1)/s$ increases on $[2,4]$.
The $D$-density is at most $\log(3)/(4\alpha r_1r_2)$.
Hence
\[
 \Delta D\le\frac{\log3}{4\alpha^2}(5\lambda/2+\beta)
                    \frac{\lambda-\mu}{\mu}\tau<5.137\tau.
\]
For an exact rational check, use $\log3<10987/10000$.
\[
\sum_{k=0}^{8}\frac{(10987/10000)^k}{k!}>3.
\]
Substitution of $\alpha,\beta,\mu,\lambda$ gives the strict $5.137\tau$ bound.
The independent Appendix~A.3 bound uses $f\le1$ on this same strip:
\[
 \Delta D\le\frac{e^{-\gamma}}{2\alpha^2}
 \left(\frac52\lambda+\beta\right)
 \log\frac{\lambda}{\mu}\,\tau
 <4.82\tau<5.137\tau.
\]
These are two bounds for the same loss. The error account retains and subtracts
only the single conservative allowance $7.160\tau$; the allocated common
error $\mathcal E_{10}$ is charged once in Section~6.
Section~6.3 assembles the endpoint and common errors.
Lemma~\ref{lem:direct-preparation} supplies preparation, with
$K_D=K_*=6(C_b+1)+2=2012\le24C_b+26$.
The coefficient identities and masks apply with the stated signed row coefficient.
\end{proof}
